\chapter{Introduction}\label{ch:intro}

\section{Motivation}

Robots can now learn manipulation skills from a few dozen human demonstrations. Early policies of this kind, such as
Action Chunking with Transformers (ACT)~\cite{zhao2023act} and Diffusion Policy~\cite{chi2023diffusion}, are small
networks trained for one task. Vision-language-action (VLA) models extend them with a pretrained vision-language
backbone that reads camera images, the robot's joint state and a written instruction, and outputs actions; examples
are \smolvla~\cite{shukor2025smolvla}, the GR00T family~\cite{nvidia2025gr00t}, \pizerofive~\cite{pi2025pi05} and
MolmoAct~2~\cite{molmoact2}. These models have from about half a billion to several billion parameters and need a GPU.

A low-cost arm, however, is driven by a small computer: a laptop or a single-board computer such as a Raspberry Pi.
The practical arrangement is \emph{remote inference}: the robot computer sends each observation over the network to a
GPU server, and the server replies with a \emph{chunk} of future actions that the robot executes at its control rate
(30~actions per second here). Because one reply carries many actions, the robot can keep moving while the next request
is in flight, which is what makes remote inference workable at all~\cite{shukor2025smolvla}.

The network is the weak point. WiFi links show latency spikes, lose messages and drop out for short periods. When a
reply is late or lost, the robot's queue of actions empties and the arm stops. The remote-inference client shipped with
LeRobot~0.7~\cite{lerobot}, the open robot-learning library this project builds on, treats a run without actions for
one second as a terminal fault and ends it, so a short outage ends the task.

Distributed Real-Time Chunking (\drtc)~\cite{vial2026drtc} is a recent scheduling method for exactly this setting: it
requests the next chunk early, estimates the latency the way TCP estimates round-trip times, and re-requests with a
fresh observation when a reply is overdue. It has not, to our knowledge, been evaluated independently against the
shipped LeRobot client on task success under controlled network faults.

\section{Problem statement}

How much do network faults lower the task success of remote inference on a low-cost robot arm, can \drtc{} recover
the lost performance, and can the robot's motion be kept smooth while it does?

\section{Objectives}

\begin{enumerate}[label=O\arabic*.]
\item Measure the effect of network faults on the task success of LeRobot~0.7's remote inference as it ships.
\item Port \drtc{} to LeRobot~0.7 and measure whether it recovers task success under the same faults.
\item Add command smoothing and measure the quality of the motion under faults.
\item Make these measurements repeatable and statistically sound: fixed plans, randomised trial order, raw logs of
      every trial, and automatic reports.
\item Train seven policies with their authors' recipes on three datasets and report their baseline performance and
      inference cost.
\end{enumerate}

The project's claim, fixed before the main experiment, is: \emph{network faults lower the task success of
client-server robot inference; \drtc{} recovers it; smoothing then improves the motion.}

\section{Contributions}

\begin{enumerate}
\item An implementation of \drtc{} inside LeRobot~0.7's remote-inference runtime that keeps all of LeRobot's own
      safety checks, with four fixes for failure modes that appeared on real hardware with long-latency, short-chunk
      models (Section~\ref{sec:drtc}).
\item Command smoothing for remote inference: a causal Butterworth filter and latency-adaptive playback
      (Section~\ref{sec:smoothing}).
\item Network fault injection at the transport and an evaluation framework: evaluation plans frozen before the first
      trial, seeded rounds, per-trial raw logs and videos, a shared results store and automatic reports with
      confidence intervals and significance tests (Sections~\ref{sec:faults} and~\ref{sec:evalsys}).
\item A controlled robot study of \R{main/n} trials (two policies, four conditions), a bench with each fault type
      alone on a replay robot, and an offline evaluation of the trained policies (Chapter~\ref{ch:results}).
\end{enumerate}

\section{Organisation of the report}

Chapter~\ref{ch:background} introduces action-chunking policies, remote inference, real-time chunking and \drtc.
Chapter~\ref{ch:system} describes the system: hardware, the baseline, our \drtc{} implementation, smoothing, fault
injection and the evaluation framework. Chapter~\ref{ch:setup} gives the datasets, policies, conditions, protocol and
metrics. Chapter~\ref{ch:results} reports the results, Chapter~\ref{ch:discussion} discusses them with their
limitations, and Chapter~\ref{ch:conclusion} concludes. The appendices list the training recipes, the full metric
definitions and how to reproduce the results.
